\section{Discussion and Limitations}\label{discussion-and-limitations}

\subsection{Demonstrated Properties}\label{demonstrated-properties}

The current artifact suite demonstrates that the control-plane protocol can run a bounded fixed-budget real-worker campaign, classify keep/discard/failed-invalid outcomes, preserve append-only provenance, recover per-trial artifact references, and expose governance metrics as paper tables and figures.

Specifically, the \passthrough{\lstinline!kdd\_main\_5trial!} real campaign demonstrates that:

\begin{itemize}
\tightlist
\item
  a \passthrough{\lstinline!claw\_style\_worker!} backed by \passthrough{\lstinline!prompt\_manager!} can execute five real ResNet-trigger training runs under governance;
\item
  two valid improving edits are kept with complete patch diffs, parsed metrics, and provenance chains;
\item
  three runtime-invalid trials are classified, recorded, and exposed as first-class audit objects without corrupting node state or leaving a pending guard;
\item
  provenance is complete for all five records (100\% rate);
\item
  the append-only ledger correctly reflects the true sequence of events.
\end{itemize}

The stress campaigns prove that scope violations (\passthrough{\lstinline!invalid\_edit\_scope!}) and no-op patches (\passthrough{\lstinline!no\_op\_patch!}) are recorded as first-class audit objects with artifact references and unchanged git state.

\subsection{Scope and Non-Claims}\label{scope-and-non-claims}

The following claims are explicitly not made:

\begin{itemize}
\tightlist
\item
  \textbf{General autonomous scientist capability.} The system governs a fixed-scope hyperparameter search loop on one node. It is not a general scientific discovery engine.
\item
  \textbf{Scientific discovery.} The net AUC gain of +0.000933 over five trials on the ResNet-trigger task is secondary evidence of real execution, not a claim of meaningful ML improvement.
\item
  \textbf{Universal optimisation algorithm.} The harness controls the experimentation process; the proposal selector is a structured round-robin, not an optimisation method.
\item
  \textbf{Backend-specific dependency.} The governance protocol does not depend on a specific coding agent, LLM, or worker implementation. Manager and worker are replaceable without changing the control-plane contract.
\item
  \textbf{Broad benchmark coverage.} All real evidence uses the ResNet-trigger node. Results may not generalise to other ML tasks without additional evaluation.
\item
  \textbf{Full three-way lifecycle diversity.} The current main campaign includes \passthrough{\lstinline!kept!} and \passthrough{\lstinline!failed\_invalid!} outcomes but no \passthrough{\lstinline!discarded!} valid-but-worse trial. A longer campaign or a valid non-improving proposal would complete the lifecycle demonstration.
\item
  \textbf{Rationale memory superiority.} The 10-trial LangGraph ablation produces a non-monotonic result: summary memory achieves the lowest repeated-bad rate (0.50) while none and rationale are tied (0.80). We do not claim that rationale memory outperforms summary memory; if anything, the evidence suggests the opposite, though a single 10-trial run per arm is insufficient to draw a firm conclusion.
\end{itemize}

\subsection{Memory Ablation Analysis}\label{memory-ablation-analysis}

The pre-stated hypothesis --- repeated\_bad\_rate(none) \textgreater{} repeated\_bad\_rate(summary) \textgreater{} repeated\_bad\_rate(rationale) --- is not confirmed in either of the two ablation configurations that were run.

\textbf{\passthrough{\lstinline!prompt\_manager!} arm (design-constrained negative).} All three arms produce identical outcomes (2 kept, 3 failed-invalid, repeated\_bad\_rate = 0.40, best val\_auc = 0.774711). The mechanism is in Section 4: the deterministic round-robin makes memory structurally irrelevant. Memory injection and repeated-bad detection work correctly; the limitation is the proposal generator.

\textbf{10-trial \passthrough{\lstinline!prompt\_manager!} extension (partial positive).} Extending the same structured manager to 10 trials changes the conclusion. The no-memory arm records repeated\_bad\_rate = 0.60, while both memory arms record repeated\_bad\_rate = 0.40 and reach a higher best validation AUC. This supports the weaker claim that memory can reduce repeated poor proposals under a longer budget. It does not support the stronger hypothesis that rationale memory beats summary memory, because the two memory arms are tied.

\textbf{\passthrough{\lstinline!langgraph\_manager!} arm (stochastic LLM, genuinely behavioural).} The v2 LangGraph campaigns (\passthrough{\lstinline!lg\_ablation2\_*!}, budget=10, temperature=0.7) use a deterministic-patch bridge that eliminates the external coding-agent dependency, producing zero \passthrough{\lstinline!edit\_failed!} outcomes across 30 trials. Memory injection is verified per-trial: all 30 SHA-256 context hashes are distinct (10 per arm), confirming that memory accumulates correctly within each arm and differs across modes. The result is non-monotonic: the \passthrough{\lstinline!append\_only\_summary!} arm achieves both the lowest repeated-bad rate (0.50 vs.\ 0.80 for both other arms) and the highest best AUC (0.8212 vs.\ 0.774711). The \passthrough{\lstinline!none!} and \passthrough{\lstinline!append\_only\_summary\_with\_rationale!} arms are tied. The \passthrough{\lstinline!edit\_failed!} failure mode from the original five-trial runs (\passthrough{\lstinline!lg\_ablation\_*!}) is a positive governance finding that remains valid: the control plane correctly refused to keep 15/15 trials where training ran against unmodified code, demonstrating fail-safe behaviour under a novel failure condition. That historical evidence is preserved in the ledgers; the v2 campaign supersedes it for the memory-effect measurement.

This is a layered finding rather than a simple win. It correctly characterises separate limits: (1) a five-trial deterministic run was too short and too constrained to expose a memory effect; (2) a 10-trial deterministic run shows memory reducing repeated-bad rate but not a rationale-over-summary advantage; (3) the 10-trial stochastic run shows a non-monotonic result --- summary memory outperforms both no memory and rationale memory --- suggesting that compact, actionable context is more effective than verbose rationale context for this LLM and task. The governance infrastructure itself --- memory injection, SHA-256 verification, repeated-bad counting, fail-safe classification --- functions correctly in all configurations.

\subsection{Why Governance Metrics Matter}\label{why-governance-metrics-matter}

Final AUC alone cannot answer whether an autonomous experiment was scientifically credible. An agent could improve a metric while repeatedly attempting invalid edits, hiding failed runs, corrupting state, or keeping decisions that cannot be audited. Governance metrics make those behaviors visible.

Acceptance rate, invalid rate, repeated-bad rate, provenance completeness, artifact capture completeness, and failure categories describe the reliability of the experimentation process around the model. This distinction is the main contribution: the harness does not replace task metrics; it makes task metrics interpretable by attaching lifecycle, validity, and provenance context to every trial.

\textbf{Governance metrics as actionable diagnostics.} The 20-trial campaign (\passthrough{\lstinline!kdd\_resnet\_scientific\_20\_trials!}) demonstrates that governance metrics are not only descriptive --- they are diagnostic. Ledger inspection of 11~\passthrough{\lstinline!proposal\_precondition\_failed!} events (all with \passthrough{\lstinline!wall\_clock\_seconds\,=\,0.0}) identified three root causes in the proposal selector without requiring any re-run: static-pool exhaustion (approx.\ 7/11), env-override shadowing (approx.\ 3/11), and stale symbol references (approx.\ 1/11). This audit-driven diagnosis is only possible because the harness records a structured \passthrough{\lstinline!extra.manager!} object --- including \passthrough{\lstinline!skipped\_candidates!} with per-candidate rejection reasons --- for every trial, including rejected ones. A naive unguarded loop would either block indefinitely (no valid proposal) or silently repeat a no-op edit; the governed control plane converts the pathological case into a first-class, inspectable ledger record. The subsequent targeted fix (expanding the candidate pool from 9 to 21 valid ResNet candidates; filtering env-pinned symbols upfront) was applied, and a follow-up 10-trial campaign (\passthrough{\lstinline!kdd\_resnet\_scientific\_10!}) confirmed the repair: \passthrough{\lstinline!proposal\_precondition\_failed!} dropped from 55\% (11/20) to 0\% (0/10), with all trials reaching real training and producing valid governance records (3~kept, 7~discarded, acceptance rate 0.30). This closed loop --- ledger evidence $\to$ targeted fix $\to$ verified re-run --- is the governance diagnostic cycle in action.

\subsection{Limitations}\label{limitations}

\textbf{Single node.} All real evidence uses the ResNet-trigger node. We cannot rule out that governance behavior reported here is specific to narrow hyperparameter editing on this node. Applying the framework across multiple nodes with holdout splits would strengthen the generalisation claim.

\textbf{Static proposal-selector candidate pool (resolved).} The structured proposal selector originally used a fixed list of 32 (symbol, target-value) pairs shared across all node types; only 9 were valid for the ResNet-trigger node with fast-search active, leading to search-space exhaustion after $\sim$10 trials. The pool was expanded to 45 entries (21 valid for ResNet-trigger) and env-pinned symbols (\passthrough{\lstinline!BATCH\_SIZE!}, \passthrough{\lstinline!EPOCHS!}) are now excluded upfront. A post-fix 10-trial campaign confirmed the repair: \passthrough{\lstinline!proposal\_precondition\_failed!} fell from 55\% to 0\%. The remaining design limitation is that the candidate list is shared globally rather than specified per node, which could require similar audits when adding new node types.

\textbf{No holdout evaluation node.} Better-Harness-style holdout evaluation nodes are not yet implemented. Without holdout nodes, it is possible for the proposal selector to overfit to the evaluation surface. This is future work.

\textbf{No discarded trial in main campaign.} The current five-trial campaign does not include a valid-but-worse trial. Full three-way lifecycle diversity (kept, discarded, failed-invalid) requires a longer campaign or a proposal that produces a valid metric below the current best.

\textbf{Memory ablation: non-monotonic, partial signal.} The 10-trial structured-manager extension shows memory reducing repeated-bad rate relative to no memory, but does not distinguish summary from rationale (both arms tied at RBR=0.40). The 10-trial LangGraph ablation (\passthrough{\lstinline!lg\_ablation2\_*!}) produces a non-monotonic result: the summary arm achieves the lowest RBR (0.50) and highest best AUC (0.8212), while none and rationale are tied (RBR=0.80). This reverses the pre-stated ordering (none $>$ summary $>$ rationale), but is based on a single 10-trial run per arm. A stronger test would replicate across multiple seeds and larger budgets, or use a task with a wider editable surface where different memory modes produce structurally distinct proposal trajectories.

\textbf{Stress artifact reproducibility.} Stress trial artifact files are generated deterministically by the stress workers and are byte-reproducible from the worker source code. The corresponding ledger records are unchanged and continue to serve as the authoritative audit objects.

\textbf{Task-metric noise.} Five baseline seeds produce validation AUCs from 0.771511 to 0.810711, with mean 0.784755 and bootstrap 95\% CI {[}0.774809, 0.798329{]}. This variation is larger than the primary five-trial campaign's +0.000933 gain, so task-metric improvement should be treated as secondary evidence that the loop runs, not as an optimisation claim.

\textbf{Reproducibility and worker dependencies.} The \passthrough{\lstinline!claw\_style\_worker!} has two execution paths. The deterministic patch path --- gated on \passthrough{\lstinline!proposal.extra["deterministic\_patch"]!} and \passthrough{\lstinline!proposal.extra["structured\_edit"]!} --- applies constant-value edits directly and requires no external coding agent. The free-form path translates LLM proposals into patches by invoking a live AI coding agent (Claude Code or a claw-harness agent). The \passthrough{\lstinline!edit\_failed!} failure mode (Section 5) arises when the free-form path is taken without an agent present: training runs on the unmodified source, and the control plane correctly refuses to keep the result.

The \passthrough{\lstinline!LangGraphManager!} deterministic patch bridge (\passthrough{\lstinline!\_extract\_structured\_edit!}, implemented in the harness) eliminates this dependency for proposals that include structured \passthrough{\lstinline!param!}/\passthrough{\lstinline!old\_value!}/\passthrough{\lstinline!new\_value!} fields. Any replication attempt using \passthrough{\lstinline!LangGraphManager!} with a model that produces structured proposals will route through the deterministic path without requiring a coding agent.

For fully dependency-free replication --- no Ollama, no coding agent --- the \passthrough{\lstinline!LocalWorker!} (\passthrough{\lstinline!src/autoresearch/worker/local\_worker.py!}) parses ``Change PARAM from X to Y'' directives from any proposal text and applies them directly. This worker exercises the full control-plane lifecycle, including scope validation, metric parsing, and ledger recording, without external tool dependencies. Reviewers wishing to verify the governance protocol can run a three-trial dry campaign with \passthrough{\lstinline!LocalWorker!} or \passthrough{\lstinline!DryRunWorker!} using only standard Python dependencies.

\subsection{Generalisation Path}\label{generalisation-path}

The NodeSpec YAML pattern generalises the harness to new ML experiments without code changes; each spec is a harness template for a class of experiments. A stronger follow-up evaluation would:

\begin{itemize}
\tightlist
\item
  add multiple node specs, including a holdout node not used for development;
\item
  replicate the memory ablation across multiple seeds to assess whether the non-monotonic summary $<$ rationale ordering observed in \passthrough{\lstinline!lg\_ablation2\_*!} is stable, and extend to a wider editable surface where different memory modes may diverge more strongly;
\item
  compare multiple manager backends (baseline, prompt, LangGraph) under equivalent budgets and worker modes;
\item
  report bootstrap confidence intervals for governance metrics across repeated seeds.
\end{itemize}
